% ==========================================
% BAB VI EVALUASI
% ==========================================
\cleardoublepage
\chapter{EVALUASI SISTEM}
\label{chap:evaluasi-sistem}

Bab ini membahas evaluasi terhadap \textit{frontend dashboard} pemantauan \textit{Water Treatment Plant} (WTP) 
yang telah diimplementasikan. Evaluasi dilakukan melalui dua tahap pengujian, yaitu pengujian fungsional 
(\textit{black-box testing}) untuk memverifikasi kesesuaian perilaku sistem terhadap spesifikasi, dan pengujian 
usabilitas menggunakan \textit{System Usability Scale} (SUS) untuk menilai kemudahan penggunaan antarmuka. 
Hasil kedua pengujian tersebut selanjutnya dipetakan terhadap kebutuhan fungsional yang telah didefinisikan 
pada Bab~\ref{chap:analisis-masalah}, dan diakhiri dengan identifikasi batasan sistem.

\section{Rancangan Pengujian}
Pengujian dirancang untuk menjawab kriteria keberhasilan yang telah ditetapkan pada Bab~\ref{chap:pendahuluan}, 
yaitu seluruh kebutuhan fungsional \textit{frontend} dapat diverifikasi melalui pengujian fungsional. Ruang lingkup 
pengujian dibatasi pada perilaku antarmuka dan alur interaksi operator. Pengujian ini tidak mencakup validasi akurasi 
kalibrasi sensor, efektivitas proses kimia, kebenaran algoritma kendali dosis pada \textit{backend}, maupun keandalan 
jaringan pada kondisi lapangan, sebagaimana telah dinyatakan dalam batasan masalah.

Pengujian fungsional dilakukan dengan metode \textit{black-box}, yaitu dengan mengeksekusi setiap skenario secara 
langsung melalui antarmuka web tanpa meninjau struktur kode di baliknya. Suatu skenario dinyatakan berhasil 
(\textit{passed}) apabila antarmuka menampilkan data dari \textit{backend} dengan benar, menerima masukan pengguna 
sesuai kebutuhan, dan memberikan respons yang sesuai terhadap aksi operator.

\section{Hasil Pengujian Fungsional}
Pengujian fungsional dilakukan terhadap 20 skenario yang mencakup autentikasi, pemantauan \textit{dashboard}, 
analitik sensor, peninjauan riwayat, ekspor laporan, penanganan alarm, pengaturan \textit{offset} pompa, serta 
penanganan masukan tidak valid. Ringkasan skenario dan hasil pengujian ditunjukkan pada Tabel~\ref{tab:blackbox}.

\begin{longtable}{p{1.2cm}p{2.6cm}p{4.2cm}p{4.2cm}p{1.4cm}}
\caption{Skenario dan hasil pengujian fungsional \textit{frontend}} \label{tab:blackbox} \\
\toprule
\textbf{ID} & \textbf{Use Case} & \textbf{Skenario Pengujian} & \textbf{Hasil yang Diharapkan} & \textbf{Status} \\
\midrule
\endfirsthead

\multicolumn{5}{c}{\tablename\ \thetable\ Skenario dan hasil pengujian fungsional \textit{frontend} (lanjutan)} \\
\toprule
\textbf{ID} & \textbf{Use Case} & \textbf{Skenario Pengujian} & \textbf{Hasil yang Diharapkan} & \textbf{Status} \\
\midrule
\endhead

\midrule
\multicolumn{5}{r}{\textit{Bersambung ke halaman berikutnya.}} \\
\endfoot

\bottomrule
\endlastfoot

FE-01 & Login & Operator memasukkan kredensial akun yang valid pada halaman login. & Sistem mengautentikasi pengguna dan mengarahkan operator ke halaman \textit{dashboard}. & \textit{Passed} \\
FE-02 & Logout & Operator menekan tombol \textit{logout}. & Sesi pengguna berakhir dan operator diarahkan kembali ke halaman login. & \textit{Passed} \\
FE-03 & Melihat \textit{Dashboard} & Operator membuka halaman \textit{dashboard} pemantauan. & \textit{Dashboard} menampilkan ringkasan kondisi sistem, status kepatuhan parameter, level tangki bahan kimia, dan status pompa dosis. & \textit{Passed} \\
FE-04 & Melihat \textit{Dashboard} & Operator melihat informasi parameter kualitas air pada \textit{dashboard}. & Nilai pH, kekeruhan, dan TDS ditampilkan melalui komponen \texttt{ComplianceSummary} sebagai bagian dari daftar pemeriksaan kepatuhan parameter. & \textit{Passed} \\
FE-05 & Melihat \textit{Dashboard} & Operator melihat ringkasan kepatuhan pada \textit{dashboard}. & Sistem menampilkan jumlah parameter aman, total parameter, dan persentase kepatuhan. & \textit{Passed} \\
FE-06 & Melihat \textit{Dashboard} & Operator melihat visualisasi level tangki bahan kimia. & Sistem menampilkan level tangki PAC, kaporit, dan soda ash. & \textit{Passed} \\
FE-07 & Melihat \textit{Dashboard} & Operator melihat status pompa dosis pada \textit{dashboard}. & Sistem menampilkan status pompa dosis. & \textit{Passed} \\
FE-08 & Melihat \textit{Dashboard} & Sistem diuji ketika data telemetri tidak tersedia atau \textit{backend} menandai sistem \textit{offline}. & Notifikasi peringatan muncul ketika status \textit{dashboard} \textit{offline} atau data terbaru tidak tersedia. & \textit{Passed} \\
FE-09 & Memantau Sensor & Operator membuka halaman sensor. & Sistem menampilkan parameter pH, kekeruhan, dan TDS dalam bagian yang terpisah. & \textit{Passed} \\
FE-10 & Memantau Sensor & Operator melihat grafik tren pada halaman sensor. & Sistem menampilkan grafik tren untuk setiap parameter. & \textit{Passed} \\
FE-11 & Memantau Sensor & Operator melihat ringkasan statistik sensor. & Sistem menampilkan nilai minimum, maksimum, dan rata-rata. & \textit{Passed} \\
FE-12 & Melihat Riwayat Data & Operator membuka halaman riwayat dan memasukkan rentang waktu secara manual. & Sistem menampilkan data sesuai rentang waktu yang dipilih. & \textit{Passed} \\
FE-13 & Melihat Riwayat Data & Operator memilih rentang cepat 24 jam, 7 hari, atau 30 hari. & Sistem menampilkan data historis sesuai rentang waktu yang dipilih. & \textit{Passed} \\
FE-14 & Ekspor Laporan & Operator mengekspor data riwayat ke format CSV. & Sistem menghasilkan dan mengunduh laporan data riwayat dalam format CSV. & \textit{Passed} \\
FE-15 & Ekspor Laporan & Operator mengekspor data riwayat ke format XLSX. & Sistem menghasilkan dan mengunduh laporan data riwayat dalam format XLSX. & \textit{Passed} \\
FE-16 & Memantau Alarm & Operator membuka halaman alarm. & Sistem menampilkan daftar alarm dan menyediakan penyaring untuk membedakan alarm aktif dan alarm terselesaikan. & \textit{Passed} \\
FE-17 & \textit{Acknowledge} Alarm & Operator menekan tombol \textit{acknowledge} pada alarm aktif. & Sistem memproses aksi dan alarm berpindah ke daftar \textit{resolved}. & \textit{Passed} \\
FE-18 & Mengatur \textit{Offset} Pompa & Operator membuka halaman \textit{offset} pompa. & Sistem menampilkan daftar pompa beserta kalkulasi dosis dan kolom masukan \textit{offset} dalam satuan mL. & \textit{Passed} \\
FE-19 & Mengatur \textit{Offset} Pompa & Operator mengisi nilai \textit{offset} dosis berupa angka valid. & Sistem menerima masukan angka dan mengirimkannya ke \textit{backend}. & \textit{Passed} \\
FE-20 & Mengatur \textit{Offset} Pompa & Operator mengisi masukan \textit{offset} dengan nilai bukan angka. & Sistem tidak memproses pengiriman karena masukan tidak valid. & \textit{Passed} \\
\end{longtable}


Berdasarkan Tabel~\ref{tab:blackbox}, seluruh 20 skenario pengujian berhasil dieksekusi dengan status 
\textit{passed}. Ringkasan hasil pengujian berdasarkan kategori fungsionalitas ditunjukkan pada 
Tabel~\ref{tab:ringkasan-uji}.

\begin{table}[H]
\centering
\caption{Ringkasan hasil pengujian fungsional berdasarkan kategori}
\label{tab:ringkasan-uji}
\begin{tabular}{|p{6cm}|c|c|}
\hline
\textbf{Kategori} & \textbf{Jumlah Kasus} & \textbf{Hasil} \\ \hline
Autentikasi & 2 & \textit{Passed} \\ \hline
Pemantauan \textit{dashboard} & 6 & \textit{Passed} \\ \hline
Analitik sensor & 3 & \textit{Passed} \\ \hline
Peninjauan riwayat dan ekspor & 4 & \textit{Passed} \\ \hline
Penanganan alarm & 2 & \textit{Passed} \\ \hline
Interaksi \textit{offset} pompa & 3 & \textit{Passed} \\ \hline
\textbf{Total} & \textbf{20} & \textbf{\textit{Passed}} \\ \hline
\end{tabular}
\end{table}

Hasil tersebut menunjukkan bahwa \textit{frontend} tidak hanya mampu menyajikan data pemantauan, tetapi juga 
mendukung aksi operator melalui kendali antarmuka yang terproteksi, serta kemampuan sistem memperingatkan operator ketika 
data telemetri tidak mutakhir dan menolak masukan yang tidak valid sebelum dikirimkan ke \textit{backend}.

\section{Hasil Pengujian Usabilitas}
Setelah seluruh fungsionalitas terverifikasi, dilakukan evaluasi terhadap kemudahan penggunaan (\textit{usability})
antarmuka menggunakan instrumen kuesioner baku \textit{System Usability Scale} (SUS) \autocite{brooke1996sus}.
Instrumen ini terdiri atas sepuluh pernyataan, yaitu lima pernyataan bernada positif pada nomor ganjil dan lima
pernyataan bernada negatif pada nomor genap, yang masing-masing dinilai menggunakan skala Likert satu sampai lima,
dengan nilai satu menyatakan sangat tidak setuju dan nilai lima menyatakan sangat setuju. Pengujian melibatkan seorang operator WTP ITB 
Jatinangor dengan pengalaman kerja lebih dari tiga tahun sebagai responden.

\begin{table}[H]
\centering
\caption{Hasil kuesioner \textit{System Usability Scale}}
\label{tab:sus}
\renewcommand{\arraystretch}{1.3}
\begin{tabular}{|c|p{9cm}|c|c|}
\hline
\textbf{No} & \textbf{Pernyataan} & \textbf{Jenis} & \textbf{Skor} \\ \hline
1 & Saya rasa saya akan menggunakan \textit{dashboard} WTP ini secara rutin. & Positif & 5 \\ \hline
2 & Saya merasa \textit{dashboard} WTP ini terlalu rumit, padahal bisa dibuat lebih sederhana. & Negatif & 2 \\ \hline
3 & Saya merasa \textit{dashboard} WTP ini mudah digunakan. & Positif & 4 \\ \hline
4 & Saya rasa saya membutuhkan bantuan dari orang lain untuk dapat menggunakan \textit{dashboard} WTP ini. & Negatif & 3 \\ \hline
5 & Saya menemukan bahwa fungsi dan fitur di \textit{dashboard} WTP ini terintegrasi dengan baik. & Positif & 4 \\ \hline
6 & Saya merasa ada terlalu banyak hal yang tidak konsisten pada \textit{dashboard} WTP ini. & Negatif & 2 \\ \hline
7 & Saya rasa operator lain akan memahami cara menggunakan \textit{dashboard} WTP ini dengan sangat cepat. & Positif & 5 \\ \hline
8 & Saya merasa \textit{dashboard} WTP ini sangat membingungkan saat digunakan. & Negatif & 3 \\ \hline
9 & Saya merasa sangat percaya diri saat menggunakan \textit{dashboard} WTP ini. & Positif & 4 \\ \hline
10 & Saya perlu mempelajari banyak hal terlebih dahulu sebelum bisa lancar menggunakan \textit{dashboard} WTP ini. & Negatif & 2 \\ \hline
\multicolumn{3}{|r|}{\textbf{Total Kontribusi}} & \textbf{30} \\ \hline
\end{tabular}
\end{table}

Kontribusi skor untuk pernyataan bernomor ganjil dihitung dengan mengurangi skor jawaban dengan satu, sedangkan 
kontribusi untuk pernyataan bernomor genap dihitung dengan mengurangkan skor jawaban dari lima. Akumulasi 
kontribusi dari kelima item ganjil adalah 17 dan dari kelima item genap adalah 13sehingga diperoleh total 
kontribusi sebesar 30. Nilai akhir SUS diperoleh dengan mengalikan total kontribusi tersebut dengan konstanta 
2,5, sebagaimana ditunjukkan pada Persamaan berikut.

\begin{equation*}
    \text{Skor SUS} = 30 \times 2{,}5 = 75
\end{equation*}

Skor akhir SUS yang diperoleh adalah 75. Nilai tersebut melampaui target minimal sebesar 70 yang ditetapkan 
sebagai kriteria keberhasilan, sekaligus berada di atas nilai rata-rata industri yang umumnya berkisar pada 
angka 68 \autocite{sauro2011sus}. Merujuk pada \textit{adjective rating scale}, skor 75 termasuk dalam kategori 
\textit{good}, sedangkan berdasarkan \textit{acceptability range} skor tersebut termasuk dalam kategori 
\textit{acceptable} \autocite{bangor2009sus}.

Ditinjau dari pola jawaban responden, penilaian tertinggi diberikan pada aspek keinginan menggunakan sistem 
secara rutin dan kecepatan operator lain dalam memahami antarmuka (keduanya bernilai 5). Sementara itu, dua 
pernyataan negatif memperoleh skor 3, yaitu terkait kebutuhan bantuan pihak teknis dan tingkat kebingungan saat 
penggunaan. Pola ini mengindikasikan bahwa meskipun antarmuka dinilai mudah digunakan secara keseluruhan, masih 
terdapat ruang perbaikan pada aspek pemanduan pengguna, misalnya melalui penambahan pesan konfirmasi yang lebih 
eksplisit pada aksi penting maupun petunjuk penggunaan pada halaman yang bersifat teknis.

Perlu dicatat bahwa pengujian usabilitas ini hanya melibatkan satu responden, yaitu operator yang bertugas pada 
WTP ITB Jatinangor. Oleh karena itu, skor yang diperoleh merepresentasikan penilaian dari pengguna sasaran 
sistem, tetapi belum dapat digeneralisasi secara statistik. Hasil ini sebaiknya dipandang sebagai indikasi awal 
penerimaan pengguna, bukan sebagai pengukuran usabilitas yang konklusif.

\section{Evaluasi Ketercapaian Kebutuhan Fungsional}
Berdasarkan hasil pengujian pada Tabel~\ref{tab:blackbox}, seluruh kebutuhan fungsional yang didefinisikan pada 
Bab~\ref{chap:analisis-masalah} dapat diverifikasi. Pemetaan antara kebutuhan fungsional dan skenario pengujian 
yang memverifikasinya ditunjukkan pada Tabel~\ref{tab:pemetaan-kf}.

\begin{table}[H]
\centering
\caption{Pemetaan kebutuhan fungsional terhadap skenario pengujian}
\label{tab:pemetaan-kf}
\begin{tabular}{|c|p{7cm}|p{3cm}|c|}
\hline
\textbf{ID} & \textbf{Kebutuhan Fungsional} & \textbf{Skenario Uji} & \textbf{Status} \\ \hline
KF-01 & \textit{Login}, \textit{logout}, dan pembatasan akses halaman & FE-01, FE-02 & Terpenuhi \\ \hline
KF-02 & Ringkasan kepatuhan dan parameter kualitas air terkini & FE-04, FE-05 & Terpenuhi \\ \hline
KF-03 & Visualisasi level tangki dan status pompa dosis & FE-06, FE-07 & Terpenuhi \\ \hline
KF-04 & Peringatan visual saat data telemetri kedaluwarsa & FE-08 & Terpenuhi \\ \hline
KF-05 & Grafik, statistik, dan riwayat parameter kualitas air & FE-09, FE-10, FE-11 & Terpenuhi \\ \hline
KF-06 & Penyaringan rentang waktu dan ekspor CSV/XLSX & FE-12, FE-13, FE-14, FE-15 & Terpenuhi \\ \hline
KF-07 & Daftar alarm aktif dan terselesaikan serta \textit{acknowledge} & FE-16, FE-17 & Terpenuhi \\ \hline
KF-08 & Informasi dosis dan masukan \textit{offset} yang valid & FE-18, FE-19, FE-20 & Terpenuhi \\ \hline
KF-09 & Navigasi antarhalaman \textit{dashboard} & FE-03, FE-09, FE-12, FE-16, FE-18 & Terpenuhi \\ \hline
\end{tabular}
\end{table}

Dengan demikian, kriteria keberhasilan yang ditetapkan pada Bab~\ref{chap:pendahuluan} telah tercapai. Rumusan masalah 
mengenai belum terpusatnya informasi operasional dijawab melalui halaman \textit{dashboard} yang menggabungkan kepatuhan 
parameter, level tangki bahan kimia, dan status pompa ke dalam satu tampilan (KF-02, KF-03). Persoalan keterlambatan 
pengenalan kondisi yang memerlukan tindakan dijawab melalui halaman alarm dengan alur pengakuan yang terstruktur (KF-07). 
Adapun risiko pengambilan keputusan berdasarkan data yang tidak valid dijawab melalui penanganan kondisi pemuatan, 
kegagalan, dan kedaluwarsa data yang diverifikasi pada skenario FE-08 (KF-04).

\section{Batasan dan Kendala Sistem}
Meskipun seluruh skenario pengujian fungsional berhasil dilalui, terdapat sejumlah batasan yang perlu diperhatikan 
dalam menafsirkan hasil evaluasi ini:
\begin{enumerate}
    \item \textbf{Ketergantungan pada Interval \textit{Polling}:} Karena \textit{frontend} memperoleh data melalui 
    penarikan berkala (\textit{polling}), visualisasi telemetri memiliki jeda waktu (\textit{delay}) minimal sebesar 
    interval \textit{polling} dibandingkan dengan kondisi aktual di lapangan. Sistem belum mengadopsi komunikasi dua arah 
    yang persisten untuk pembaruan data secara seketika.
    \item \textbf{Ketergantungan Mutasi pada Ketersediaan \textit{Backend}:} Fitur pengakuan alarm dan pengiriman nilai 
    \textit{offset} dosis bergantung sepenuhnya pada ketersediaan layanan \textit{backend}. Apabila koneksi ke REST API 
    terputus, \textit{frontend} belum dapat menyimpan perubahan tersebut secara lokal untuk dikirimkan kembali setelah 
    koneksi pulih (\textit{offline mutation queue}).
    \item \textbf{Cakupan Pengujian Terbatas pada Lapisan Antarmuka:} Pengujian yang dilakukan memverifikasi perilaku 
    \textit{frontend} terhadap spesifikasi, tetapi tidak mengevaluasi akurasi kalibrasi sensor, efektivitas dosis kimia, 
    kebenaran algoritma kendali pada \textit{backend}, maupun keandalan jaringan pada kondisi operasional pabrik. 
    Hasil evaluasi ini karenanya harus dipahami sebagai verifikasi subsistem \textit{frontend}, bukan validasi keseluruhan 
    sistem otomasi WTP.
    \item \textbf{Jumlah Responden Pengujian Usabilitas:} Evaluasi usabilitas hanya melibatkan satu responden sehingga 
    hasilnya belum memiliki kekuatan generalisasi statistik.
\end{enumerate}